% Generated draft: requires mathematical and source review.
\documentclass[11pt]{article}
\usepackage[margin=25mm]{geometry}
\usepackage{fontspec}
\setmainfont{Latin Modern Roman}
\usepackage{amsmath,amssymb,mathtools,unicode-math}
\setmathfont{Latin Modern Math}
\usepackage{enumitem,needspace,xurl,hyperref}
\setlength{\parindent}{0pt}
\setlength{\parskip}{5pt}
\setlength{\emergencystretch}{3em}
\newcommand{\N}{\mathbb{N}}
\newcommand{\Z}{\mathbb{Z}}
\newcommand{\Q}{\mathbb{Q}}
\newcommand{\R}{\mathbb{R}}
\newcommand{\C}{\mathbb{C}}
\newcommand{\F}{\mathbb{F}}
\newcommand{\A}{\mathbb{A}}
\newcommand{\PP}{\mathbb{P}}
\newcommand{\E}{\mathbb{E}}
\newcommand{\eps}{\varepsilon}
\newcommand{\dd}{\,\mathrm{d}}
\begin{document}
\section*{UnsolvedMath ID: 30006082}
\textbf{Source problem code:} catalogue entry
\subsection{Definitions and mathematical statement}
The supplied record concerns an adjoint-free method for approximating a PDE solution operator. The method queries that operator on eigenfunctions of a prior self-adjoint operator, with the Laplace--Beltrami operator given as an example. According to the supplied source excerpt, PDE regularity estimates control the range of the adjoint, and the resulting approximation error decays at a rate determined by the eigenvalues of the prior operator. These are reported source claims; their precise hypotheses and estimates are not included in the record.

For descriptive purposes only, let \(\delta\geq 0\) denote the size of the perturbation and let \(B(\delta)\) denote the approximation-error bound referred to in the source. Neither the perturbed object, the definition of perturbation size, nor the formula for \(B\) is supplied.

\textbf{Requested conjecture.} Under the hypotheses of the source's adjoint-free approximation method, extend the asserted linear dependence of its approximation-error bound on perturbation size from small perturbations to large perturbations.

A possible schematic interpretation of an at-most-linear bound is
\[
B(\delta)\leq A+C\delta
\qquad\text{for every admissible perturbation size }\delta,
\]
where \(A,C\geq 0\) are independent of \(\delta\). This display is an interpretive template, not a recovered source inequality: the record does not specify whether this is the intended meaning of linear growth, which perturbations are admissible, or what other parameters the constants may depend on. It also does not establish that an error bound depends only on perturbation size rather than on the perturbation itself. Consequently, a precise quantified conjecture cannot be recovered without the underlying estimate and its hypotheses.
\subsection{Short English statement}
The conjecture asks whether the source's approximation-error estimate for adjoint-free operator learning retains linear dependence on perturbation size beyond the small-perturbation regime. The supplied record labels this extension open, but omits the estimate, perturbation model, norms, and admissibility conditions needed for a precise mathematical statement.
\subsection{Sources}
\begin{itemize}
\item Supplied local UnsolvedMath registry record, permanent ID \(30006082\), title \emph{Linear Error Growth for Large Operator Perturbations}, source problem code \texttt{catalogue entry}, in the Oberwolfach Reports Open Problems collection. The record gives proposed year 2024 and supplies no named proposer.
\item Supplied source citation: \emph{Deep Learning for PDE-based Inverse Problems} (2025), Oberwolfach Reports, pp.\ 2805--2900, DOI \texttt{10.4171/owr/2024/48}. \url{https://doi.org/10.4171/owr/2024/48}. Also listed among the supplied literature sources as \emph{Oberwolfach Report 48/2024}. The record describes Oberwolfach Reports as documenting mathematical workshops and their open-problem discussions.
\item Supplied literature source: Boullé et al., \emph{Operator learning without the adjoint}. \url{https://www.jmlr.org/papers/v25/24-0162.html}.
\item Literature assessment dated \texttt{2026-08-22}: the supplied curation labels the problem open, states that linear error growth is proved only for small perturbations and conjectured for large perturbations, and reports that current literature supplies no general proof of the extension. This is supplied dated curation, not an independently verified literature assessment.
\item The registry labels its statement \texttt{corrected\_verified} and says that it was reconstructed and checked against the official Oberwolfach report. This is supplied verification metadata; no source visit or independent verification is claimed here.
\end{itemize}
\end{document}
